\begin{lemma}[Scheide defective no-cross property (4)]
For every defective color $\delta$ in a defective-color setup, no
$\delta$-colored edge identity joins
$U_{\gamma(\delta)}$ to $Z_{\gamma(\delta)}$.
\end{lemma}
\begin{proof}
Fix a defective color $\delta$, write
\[
 X=\noderef{TreeSequenceVertexSet}(T),\qquad
 i=\gammaIndex(\delta),\qquad \gamma=\gamma_i,
\]
and put
\[
 U=\{u^1_\gamma,u^2_\gamma\},\qquad Z=X\setminus U.
\]
We first unpack all data that will be used.  The definition
\noderef{ScheideDefectiveSetup} supplies a color $\alpha_0$ missing at the
first base vertex $y_0$, with $\alpha_0\ne\gamma_r$ for every block index
$r$.  It supplies a full simple $(\alpha_0,\delta)$ incidence path
\[
 P=P_{y_0}(\alpha_0,\delta,\varphi)
\]
from $y_0$ to a vertex $z_\delta\notin X$.  On this recorded path,
$v^0=v^0_\delta$ is the last vertex in $X$, and the next two incidence
nodes are the edge-node $f^0=f^0_\delta$ and the vertex-node
$u^0=u^0_\delta$.  In particular,
\[
 \varphi(f^0)=\delta,\qquad
 \operatorname{ends}(f^0)=\{v^0,u^0\},\qquad u^0\notin X.       \tag{1}
\]
The same setup says that $v^0$ is a Scheide exit vertex and that
$\gamma$ is missing at $v^0$.  Its full-path clause is the predicate
\noderef{ScheideDefectiveChainPath}; consequently $P$ is simple and its
support is exactly
\[
 \noderef{TwoColorComponent}
   (G,\varphi,\alpha_0,\delta,\operatorname{inl}y_0).           \tag{2}
\]

Unpack the extension \noderef{ScheideReducedConfiguration} to the
underlying \noderef{ScheideStructuralConfiguration}.  It gives the exact
ordered decomposition of $T$ into the $p\geq2$ base vertices and edges,
followed by the two vertices and two edges of each block.  It also gives
the injectivity of the block colors, a base vertex missing each block
color, the block attachment and color equations, the avoidance of all
block colors by every colored nonfirst base edge, and a globally maximum
triple.  Expanding \noderef{GlobalMaximumTriple}, denote its
\noderef{KTriple} component by $K$; it says in particular that $\varphi$
is proper in the sense of \noderef{IsProperEdgeColoring} and satisfies
\noderef{IsOneEdgeUncolored}$(\varphi,e_0)$.  It also supplies
\noderef{TashkinovTree}$(G,e_0,\varphi,T)$ and
\[
 |X|=\noderef{TashkinovOrder}(G).                               \tag{3}
\]
By \noderef{ScheideGlobalMaximumOrder}, the same $X$ is elementary and
closed in the senses of \noderef{ColoredVertexSetElementary} and
\noderef{ColoredVertexSetClosed}.

We shall need the precise location of $v^0$, so we derive it rather than
assuming it.  Membership of $v^0$ in
\noderef{ScheideExitVertices}, supplied by the setup, unfolds through
\noderef{ScheideExitVertex} to give $v^0\in X$.  The exit-covering field
of the structural configuration supplies an index $r$ for which
$\gamma_r$ is missing at $v^0$, while its block-missing field supplies a
base vertex $y_q$ at which the same $\gamma_r$ is missing.  The position
equations and \noderef{TreeSequenceVertexSet} put $y_q$ in $X$.
Elementariness of $X$ forces $v^0=y_q$.  Thus
\[
 v^0\text{ belongs to the base }Y.                             \tag{4}
\]
The position equations and vertex injectivity in
\noderef{TreeSequence} separate every base position from every block
position.  Hence the base vertices $v^0$ and $y_0$ are different from
both vertices of $U$, while both vertices of $U$ belong to $X$ and are
distinct.  In particular,
\[
 v^0,y_0\in Z.                                                 \tag{5}
\]

We next record all needed color inequalities.  We already have
$\alpha_0\ne\gamma$.  The membership of $\delta$ in
\noderef{DefectiveColors}, together with the last clause of
\noderef{ScheideGlobalMaximumBoundaryFacts}, gives at least three
identities in
\[
 \noderef{BoundaryColorEdges}(G,\varphi,X,\delta).
\]
Choose one of them.  Closedness of $X$, unfolded via
\noderef{ColoredVertexSetClosed}, says that the color of this boundary
edge is missing at no vertex of $X$.  Because
\noderef{MissingColors} is the complement of
\noderef{PresentColors}, $\delta$ is present at every vertex of $X$.
It is therefore distinct from $\alpha_0$, which is missing at $y_0$, and
from $\gamma$, which is missing at $v^0$.  Thus
\[
 \alpha_0,\gamma,\delta\quad\hbox{are pairwise distinct}.       \tag{6}
\]

Suppose now, for a contradiction, that a $\delta$-colored edge identity
$g_1$ joins $U$ to $Z$.  By the definition \noderef{Incident}, its two
incidences and the loopless endpoint condition in
\noderef{LooplessMultigraph} identify its endpoint pair as
$\{a,z_1\}$ for some $a\in U$ and $z_1\in Z$.  Relabel the two members
of $U$, and their corresponding two deleted block pairs, locally as
$u^1=a$ and $u^2$; both block pairs will be deleted later, so this
relabeling changes no retained ordered datum.  Since
$\alpha_0$ is missing at $y_0\in Z$ and $\alpha_0\ne\gamma$,
\noderef{ScheideStructuralInternalPair} supplies an individual edge
identity $g_2$ with
\[
 \varphi(g_2)=\alpha_0,\qquad
 \operatorname{ends}(g_2)=\{u^1,u^2\}.                         \tag{7}
\]

We choose the third identity with the required endpoint information.
Equation (1), (4), and (5), together with the definitions
\noderef{BoundaryEdges} and \noderef{BoundaryColorEdges}, show that
\[
 f^0\in\noderef{BoundaryColorEdges}(G,\varphi,X,\delta)
 \quad\hbox{and its endpoint in }X\hbox{ lies in }Z.           \tag{8}
\]
At $u^1$, the identity $g_1$ already occupies the $\delta$-incidence.
No different $\delta$-edge can be incident there: two such identities
would be adjacent by \noderef{AdjacentEdges} and \noderef{Incident},
contrary to properness.  Properness likewise permits at most one
$\delta$-edge incident with $u^2$.  A boundary identity has exactly one
endpoint in $X$ by \noderef{BoundaryEdges}; hence at most one member of
$\noderef{BoundaryColorEdges}(G,\varphi,X,\delta)$ has its endpoint in
$U$.  If every member other than $f^0$ had its endpoint in $U$, that
boundary-color set would have cardinality at most two.  This contradicts
its cardinality at least three.  We may therefore choose an identity
$g_3\ne f^0$ and vertices $z_3,w_3$ such that
\[
 g_3\in\noderef{BoundaryColorEdges}(G,\varphi,X,\delta),\qquad
 z_3\in Z,\quad w_3\notin X,\quad
 \operatorname{ends}(g_3)=\{z_3,w_3\}.                         \tag{9}
\]
This argument counts edge identities, so it remains valid for parallel
edges.  Notice also that $g_1$ and $g_2$ are internal to $X$, whereas
$g_3$ and $f^0$ are boundary identities.  Together with (6)--(9), this
shows already that $g_1,g_2,g_3$ are pairwise distinct.

Both $\alpha_0$ and $\gamma$ are missing somewhere on $X$, at $y_0$ and
$v^0$ respectively.  By \noderef{TreeMissingColors}, they belong to the
tree-missing set.  The missing-color clause of
\noderef{ScheideGlobalMaximumBoundaryFacts} therefore gives
\[
 \noderef{BoundaryColorEdges}(G,\varphi,X,\alpha_0)=\varnothing,
 \qquad
 \noderef{BoundaryColorEdges}(G,\varphi,X,\gamma)=\varnothing. \tag{10}
\]
Define the first recoloring exactly by
\[
 \varphi_1=\noderef{OutsideTreeColorSwap}
   (T,\varphi,\alpha_0,\gamma).
\]
The definition fixes an identity satisfying
\noderef{TreeIncidentEdge}$(T,f)$ and applies the transposition of the
distinct colors $\alpha_0,\gamma$ to every other edge identity.  In
particular, every edge incident with a vertex of $X$ is fixed.

We verify properness of $\varphi_1$ edge by edge.  Let $a,b$ be distinct
adjacent colored identities under $\varphi_1$.  By
\noderef{AdjacentEdges}, some vertex $x$ is incident with both.  If both
identities are incident with $X$, both values are fixed.  If neither is,
the same injective transposition is applied to both old values.  Equality
of their new colors would in either case imply equality of their old
colors, contradicting properness of $\varphi$.  Finally suppose exactly
one identity, say $a$, is incident with $X$.  The common vertex $x$ must
lie outside $X$, since otherwise $b$ would also be incident with $X$.
Thus $a$ is a boundary edge in the sense of \noderef{BoundaryEdges}.
Equation (10) says that its fixed color is neither $\alpha_0$ nor
$\gamma$.  Transposing those two colors on $b$ can equal that fixed value
only if $b$ already had the same value, again contradicting properness of
$\varphi$.  This proves
\[
 \noderef{IsProperEdgeColoring}(G,\varphi_1).                  \tag{11}
\]
The unique-uncolored condition is equally explicit.  The $K$-triple gives
$\varphi(e_0)=\mathsf{none}$ and a colored value on every other identity.
The outside swap fixes $\mathsf{none}$ and maps every colored value to a
colored value.  Applying this to $e_0$ and then to each $f\ne e_0$ proves
\[
 \noderef{IsOneEdgeUncolored}(\varphi_1,e_0).                  \tag{12}
\]

Let $R=v^0Pz_\delta$ be the terminal subwalk of the recorded path starting
at its recorded occurrence of $v^0$.  The successive-index fields in
\noderef{ScheideDefectiveSetup} say that its first edge-node is precisely
$f^0$ and its following vertex-node is $u^0$.  The last-tree-vertex field
says that every later vertex-node of $R$ lies outside $X$.  Hence
\[
 V(R)\cap X=\{v^0\}.                                         \tag{13}
\]
The first edge $f^0$ is incident with $X$, so the outside swap fixes it;
by (1) it retains color $\delta$.  Every later edge of $R$ has both ends
outside $X$.  Its old color is $\alpha_0$ or $\delta$ by
\noderef{TwoColorIncidenceGraph}.  Along the simple proper old path,
successive edge identities have different colors, so these colors
alternate.  The outside swap changes every such outside
$\alpha_0$-edge to $\gamma$ and fixes every $\delta$-edge, using (6).
Thus $R$ is a $(\gamma,\delta)$ incidence walk under $\varphi_1$.

We now prove that it is the entire new component, not merely a walk inside
that component.  First consider its initial vertex $v^0$.  The color
$\gamma$ is missing there and the outside swap fixes every incident edge.
The identity $f^0$ is an incident $\delta$-edge; properness makes it the
unique such edge.  If $v^0=y_0$, then $\alpha_0$ is also missing there,
and $v^0$ plainly has only the $f^0$ adjacency in the new two-color graph.
If $v^0\ne y_0$, simplicity of $P$ gives a predecessor edge at $v^0$,
distinct from $f^0$.  Since $f^0$ has old color $\delta$, properness and
the two-color path condition force that predecessor to have old color
$\alpha_0$.  It is incident with $X$, so it remains $\alpha_0$ under
$\varphi_1$ and is excluded from the new $(\gamma,\delta)$ graph by (6).
Thus in both cases $v^0$ has degree one in the new component, through
$f^0$.

At every later internal vertex-node of $R$, the two path edges have the
two different colors $\gamma$ and $\delta$.  Properness (11) permits no
additional incident edge of either color.  At every internal edge-node,
the bipartite definition \noderef{TwoColorIncidenceGraph} permits only
its two endpoint incidences, both already used by $R$.  Hence no branch
leaves $R$ at an internal node.

It remains to check the terminal vertex $z_\delta$.  In the old
$(\alpha_0,\delta)$ incidence graph, $y_0$ has degree one: $\alpha_0$ is
missing there, while $\delta$ is present there and properness makes its
incident $\delta$-edge unique.  The theorem
\noderef{TwoColorIncidenceDegreeAtMostTwo} bounds every node degree in
that old graph by two.  By (2), the simple path $P$ has support equal to
the whole finite connected component.  Starting at the degree-one node
$y_0$, follow the unique unused adjacency.  At each later node the
degree-two bound leaves at most one continuation, so no branch is
possible.  Finiteness forces the traversal to stop unless it revisits a
node.  A first revisit would create a cycle; the first node where the
route from $y_0$ met that cycle would have the entering edge and two cycle
edges, contradicting the degree-two bound, while $y_0$ itself cannot lie
on a cycle because its degree is one.  Connectedness and (2) therefore
make $P$ the whole component path and make its other endpoint
$z_\delta$ a degree-one node.

Every new $\gamma$- or $\delta$-edge incident with $z_\delta$ corresponds
to an old $\alpha_0$- or $\delta$-edge there.  Indeed, on an edge having
both ends outside $X$ this is exactly the $\alpha_0,\gamma$
transposition; on a boundary edge the value is fixed, and (10) excludes
both $\alpha_0$ and $\gamma$.  The old degree-one conclusion thus gives
degree one at $z_\delta$ in the new graph.  We have checked the first
edge, the possible excluded predecessor at $v^0$, every internal vertex-
and edge-node, and the terminal endpoint.  Therefore $R$ is connected and
no incidence adjacency leaves its support.  By
\noderef{TwoColorComponent},
\[
 \operatorname{supp}(R)=
 \noderef{TwoColorComponent}
   (G,\varphi_1,\gamma,\delta,\operatorname{inl}v^0).           \tag{14}
\]

Define the second recoloring by switching exactly this full component:
\[
 \varphi_2=\noderef{TwoColorComponentSwap}
   (G,\varphi_1,\gamma,\delta,\operatorname{inl}v^0).
\]
The colors are distinct by (6), and (11) is proper, so
\noderef{TwoColorComponentSwapProper} makes $\varphi_2$ proper.  The
definition \noderef{TwoColorComponentSwap} fixes an uncolored value and
maps each colored value to a colored value.  Applying it edge by edge to
(12) proves
\[
 \noderef{IsOneEdgeUncolored}(\varphi_2,e_0).                  \tag{15}
\]

We determine the exact missing-color effect at $v^0$.  The first swap
fixed every edge incident there, so its present and missing sets under
$\varphi_1$ equal those under $\varphi$.  By (13)--(14), the second swap
contains exactly one edge at $v^0$, namely $f^0$; before the switch it has
color $\delta$, and no incident edge has color $\gamma$.  It changes
$f^0$ to $\gamma$ and fixes every other incident identity.  Consequently
\[
 \noderef{MissingColors}(G,\varphi_2,v^0)
  =\operatorname{insert}\delta
     \bigl(\noderef{MissingColors}(G,\varphi,v^0)
       \setminus\{\gamma\}\bigr).                             \tag{16}
\]
In particular, $\delta$ is missing at $v^0$.  If $x\in X$ and
$x\ne v^0$, then (13)--(14) put its vertex-node outside the switched
component.  An incident changed edge-node would connect it to the
component in \noderef{TwoColorIncidenceGraph}, which is impossible.
Thus every other vertex of $X$ has exactly its old present and missing
sets.  Equation (16) is the only missing-status change on $X$.

We also locate all switched edge identities.  Equations (13)--(14) say
that the switched component meets $X$ only at the base vertex $v^0$ and
contains only $f^0$ among the edges incident there.  Hence it contains no
edge identity having two endpoints in $X$.  In particular it contains no
retained tree edge.  The same conclusion can be checked from colors:
\noderef{ScheideGlobalMaximumUnused}, read through
\noderef{TreeUsedColors}, says that the defective color $\delta$ is unused
on the colored edges of $T$; after deleting the $i$-block, the structural
base-avoidance, block-color, and color-injectivity fields say that
$\gamma$ is also unused on every retained colored tree edge.  The first
edge is $e_0$ and is uncolored by $K$.  Thus no retained edge-node can
belong to a $(\gamma,\delta)$ component.

The first swap fixes $g_1,g_2,g_3$, because each is incident with $X$.
None lies in the component (14).  If $g_1$ did, its incidence with
$u^1\in X$ would put that vertex-node in the same component, contradicting
$u^1\ne v^0$ and (13).  The identity $g_2$ has color $\alpha_0$, outside
the pair $\{\gamma,\delta\}$ by (6), whereas every edge-node in the
nontrivial component (14) has one of its two defining colors; this excludes
$g_2$.  For $g_3$, first note that its tree endpoint $z_3$ is not $v^0$.
Otherwise $g_3\ne f^0$ would give two distinct
incident $\delta$-edges at $v^0$ under the proper coloring $\varphi_1$:
both boundary identities are fixed by the outside swap.  Thus
$z_3\ne v^0$, and its incidence excludes $g_3$ from (14).  Both
recolorings therefore leave these three exact identities unchanged, and
(7), (9) give
\[
 \varphi_2(g_1)=\delta,\qquad
 \varphi_2(g_2)=\alpha_0,\qquad
 \varphi_2(g_3)=\delta.                                      \tag{17}
\]

We now construct the ordered sequence used for the contradiction and
verify every clause rather than invoking an unspecified preservation
principle.  The structural position equations write the old sequence as
the base segment followed, in increasing block index, by the pairs
\[
 (f^1_{\gamma_r},u^1_{\gamma_r}),
 (f^2_{\gamma_r},u^2_{\gamma_r}).
\]
Let $T^-$ retain the whole base and every block pair with $r\ne i$, in
their original relative order, and delete exactly
\[
 (f^1_\gamma,u^1),\qquad(f^2_\gamma,u^2).                     \tag{18}
\]
Because $p\ge2$, a base edge remains, so the new sequence has positive
length and its first edge is still $e_0$.  Its vertices are precisely
$Z$.  Vertex injectivity and edge injectivity are inherited from
\noderef{TreeSequence}.  Every base attachment is unchanged.  Every
retained block edge has, by the structural attachment field, endpoints
consisting of its block vertex and a base vertex; that base vertex occurs
earlier in $T^-$.  Thus every attachment remains valid after the explicit
reindexing, and $T^-$ is a tree-sequence with vertex set $Z$.

Every retained edge identity has the same color under $\varphi_2$ as
under $\varphi$: the outside swap fixed it because it is incident with
$X$, and the component swap contains no retained tree edge.  We verify
the Tashkinov witnesses by the two possible locations of a retained edge.
For a nonfirst base edge, use its old witness from the
\noderef{TashkinovTree} clause of the global maximum.  Every earlier
vertex lies in the base and is retained.  The structural base-avoidance
field says that this edge's color is not $\gamma$; hence even if its old
witness is $v^0$, the only possible loss in (16) does not remove that
color.  For an edge in a retained block $r\ne i$, use the base vertex
missing $\gamma_r$ supplied by the structural block-missing field.  That
base vertex precedes the block after reindexing.  Injectivity of the block
colors gives $\gamma_r\ne\gamma$, so (16) again preserves the missing
witness, including when it is $v^0$.  These cases exhaust all nonfirst
edges of $T^-$.  Therefore
\[
 \noderef{TashkinovTree}(G,e_0,\varphi_2,T^-).                 \tag{19}
\]

Append the exact three pairs
\[
 T'=T^-\mathbin{\Vert}(g_1,u^1,g_2,u^2,g_3,w_3)               \tag{20}
\]
in that order.  We check the \noderef{TreeSequence} fields.  The first
new edge $g_1$ has endpoints $z_1\in Z$ and $u^1$, so the retained
occurrence of $z_1$ is its earlier attachment and $u^1$ is its new vertex.
The edge $g_2$ has endpoints $u^1,u^2$, so the just-appended $u^1$ is its
earlier attachment and $u^2$ is its new vertex.  The edge $g_3$ has
endpoints $z_3\in Z$ and $w_3\notin X$, so the retained occurrence of
$z_3$ is its earlier attachment and $w_3$ is its new vertex.

The new vertices are distinct from all earlier ones: $u^1,u^2$ are the
two distinct members deleted from $X$, and $w_3\notin X$.  The new edge
identities are pairwise distinct as already observed.  They are also
distinct from every retained tree edge.  The $\delta$-identities $g_1$
and $g_3$ cannot be retained colored edges by
\noderef{ScheideGlobalMaximumUnused}, and neither is $e_0$ because both
are colored while $e_0$ is uncolored.  The identity $g_2$ cannot be a
base edge or an edge of another block because its exact endpoints are the
two selected block vertices, which the structural position equations and
old vertex injectivity separate from all base vertices and all other block
vertices.  It cannot be either removed block edge because its color is
$\alpha_0\ne\gamma$.  Thus edge injectivity holds in (20).  All old
attachments and the first-edge equation persist from $T^-$.  This proves
every tree-sequence clause for $T'$.

It remains to check the three new Tashkinov witnesses.  Equations
(16)--(17) say that $g_1$ has color $\delta$ missing at the earlier
retained base vertex $v^0\in Z$.  The edge $g_2$ has color $\alpha_0$.
The first outside swap fixes every edge incident with $y_0$, and the
second swap changes at $v^0$ only the $\gamma,\delta$ statuses in (16).
Since $\alpha_0$ is distinct from both colors by (6), $\alpha_0$ remains
missing at the earlier retained vertex $y_0\in Z$.  Finally $g_3$ again
has color $\delta$ missing at the earlier vertex $v^0$.  Together with
(19), these are witnesses for every nonfirst edge of (20), so
\[
 \noderef{TashkinovTree}(G,e_0,\varphi_2,T').                 \tag{21}
\]

We rebuild the $k$-triple for the final coloring.  Criticality of $G$, the
chromatic-index equality, and the maximum-degree inequality are the three
coloring-independent clauses of the original \noderef{KTriple} $K$.
The remaining clauses are precisely the properness of $\varphi_2$ and
(15).  Hence $\noderef{KTriple}(G,k,e_0,\varphi_2)$ holds.  Its
chromatic-index equality places $T'$ among the trees maximized in
\noderef{TashkinovOrder}, so (21) gives
\[
 |V(T')|\le \noderef{TashkinovOrder}(G).                       \tag{22}
\]
On the other hand, (18) deletes exactly the two distinct vertices of $U$
and (20) appends those same two vertices and the new vertex
$w_3\notin X$.  Therefore
\[
 |V(T')|=|X|-2+3=|X|+1.                                      \tag{23}
\]
Equations (3), (22), and (23) give $|X|+1\le|X|$, a contradiction.
Thus the assumed $\delta$-colored identity joining $U$ to $Z$ cannot
exist.  Since $\delta$ and the alleged identity were arbitrary, the
required no-cross property holds for every defective color.
\end{proof}
